\scanpage{153}
\mspage{11}
2. \inserted{\unsure{les}}\uline{Catégories loc. filtrantes} \uline{sont} \uline{les catégories loc. \deleted{\unsure{filtrantes}} \inserted{tot. \unsure{0-connexes}}} ?

Soit $\mathscr T$ un topos. On dit que $\mathscr T$ est \uline{irréductible} \deleted{\unsure{s'il est dans}} si l'ens. ordonné \inserted{$\mathrm{Ouv}(\mathscr T)$,} des sous-\deleted{\unsure{b}}objets \deleted{\unsure{ouverts}} \inserted{de $e_{\mathscr T}$} \deleted{\unsure{de $\mathscr T$}} a la propriété suivante :
a) $\mathrm{Ouv}(\mathscr T)$ n'est pas réduit à un élément, i.e. $\varnothing_{\mathscr T}\to e_{\mathscr T}$ n'est pas un isomorphisme ($\mathscr T$ n'est pas identique \deleted{\unsure{au}} au sous-topos vide).
b) Si $U',U''$ sont deux objets ``non vides'' i.e. non initiaux de $\mathrm{Ouv}(\mathscr T)$, leur \unsure{produit,} intersection est ``non vide''.
Il revient au $=$ de dire que si $\mathscr T',\mathscr T''$ sont deux sous-topos fermés distincts de $\mathscr T$, alors \deleted{$\mathscr T'\vee\mathscr T''$} leur \unsure{sup} $\mathscr T'\vee\mathscr T''$ est distinct de $\mathscr T$, \inserted{ou encore que tt ouvert de $\mathscr T$ est connexe.}
% slanted
\margin{\uline{NB} Irréductible n'implique \unsure{pas} loc irréductible, pas plus que connexe n'implique loc. connexe etc}

On dit que $\mathscr T$ est \uline{loc. irréductible} si \deleted{tout objet \unsure{0-conn}}
a) $\mathscr T$ est loc connexe, \deleted{et si tout objet} (i.e. chaque faisceau sur $\mathscr T$ est somme de faisceaux 0-connexes),
et si \inserted{b)} pour tt faisceau 0-connexe \inserted{$U$} sur $\mathscr T$, le topos induit $\mathscr T_{/U}$ est irréductible \unsure{,|.} [i.e. l'ens \deleted{\unsure{ord}} ordonné $\mathrm{Ouv}(U)$ des sous-objets de $U$\unsure{,} a la propriété b) ci-dessus (a) signifiant \unsure{simpl\textsuperscript{t}} que \deleted{$U\unsure{\ne}\varnothing_{\mathscr T}$} $U\unsure{\not\simeq|\ne}\varnothing_{\mathscr T}$\unsure{\}}, i.e. que $U$ n'est pas objet initial dans $\mathscr T^{\wedge}$\unsure{\}}, ou plus \unsure{simple|simplet},] : \deleted{un sous-objet} tout sous-faisceau d'un faisceau \deleted{connexe} \unsure{est} connexe.
% margin note written sideways
\margin{$\mathscr T$ est loc irréductible ssi il a une famille génératrice $(U_\alpha)$, avec les $U_\alpha$ ``irréductibles'' (\unsure{donc} $0$-connexes -- \unsure{donc} $\mathscr T$ est loc. \deleted{\unsure{con}} connexe \unsure{?})}

Quand $\mathscr T$ est de la forme \deleted{\unsure{Top}} $\mathrm{Top}(T)$, où $T$ est un espace topologique, dire $\mathscr T$ irréductible signifie que $T$ est irréductible (car alors $\mathrm{Ouv}(\mathscr T)\simeq\mathrm{Ouv}(T)$\deleted{\unsure{\})}})\unsure{;} \deleted{et le fait \unsure{que}} \deleted{\unsure{connexe tt}} dire \inserted{\unsure{triv.}} que $T$ est loc connexe,
\deleted{et $\mathscr T$ est loc. connexe} (\deleted{\unsure{tout}} tout ouvert est somme \inserted{disj.} d'ouverts connexes), donc $\mathrm{Top}(T)$ est loc connexe. De plus, cel\unsure{a} implique aussi que $\mathrm{Top}(T)$ est loc. irréductible*. (\uline{NB} les objets 0-connexes
% slanted
\margin{expliciter la démonstration.}
% archive sheet number: 153
